% Part: sets-functions-relations
% Chapter: arithmetization
% Section: From N to Z

\documentclass[../../../include/open-logic-section]{subfiles}


\begin{document}

\olfileid{sfr}{arith}{int}
\olsection{Saka $\Nat$ menyang $\Int$}

Ana rong pangamatan dhasar:
	\begin{enumerate}
		\item Saben wilangan wutuh bisa ditulis kanthi wujud $n - m$, kanthi $n, m \in\Nat$.
		\item Informasi kang dikodeake ing ekspresi $n - m$ uga bisa dikodeake kanthi pasangan urut $\tuple{n, m}$.
	\end{enumerate}
Kita wis ngerti yen pasangan urut wilangan natural iku !!{element}s saka $\Nat^2$. Kita uga nganggep yen kita wis mangerteni $\Nat$. Mula ana usul naif adhedhasar rong pangamatan mau: \emph{ayo nganggep wilangan wutuh minangka pasangan urut wilangan natural}.

Sejatine usul iki naif banget. Mesthi wae kita kepengin $0- 2 = 4 - 6$. Nanging cetha yen $\tuple{0, 2 }\neq \tuple{4, 6}$. Dadi kita ora bisa mung nyebut $\Nat^2$ minangka himpunan wilangan wutuh.

Kanthi nggeneralisasi masalah sadurunge, kang dikarepake yaiku:
	$$a - b = c - d \text{ yen lan mung yen }a + d = c + b$$
(Cetha yen iki carane wilangan wutuh \emph{kudune} tumindak: cukup tambahna $b$ lan $d$ ing loro-lorone sisih.) Cara gampang kanggo njamin tumindak iki yaiku netepake relasi ekuivalensi antarane pasangan urut, yaiku $\Intequiv$, kaya mangkene:
$$\tuple{ a, b } \Intequiv \tuple{c, d}\text{ yen lan mung yen }a + d = c + b$$
Saiki kita kudu nuduhake yen iki relasi ekuivalensi.
\begin{prop} $\Intequiv$ iku relasi ekuivalensi.
\end{prop}
	\begin{proof}
	Kita kudu nuduhake yen $\Intequiv$ iku refleksif, simetris, lan transitif.

	\emph{Refleksif:} Cetha yen $\tuple{a, b} \Intequiv \tuple{a, b}$, amarga $a + b = b + a$.

	\emph{Simetris:} Upamakna $\tuple{a, b} \Intequiv \tuple{c, d}$, mula $a + d = c + b$. Banjur $c + b = a + d$, saengga $\tuple{c, d} \Intequiv \tuple{a, b}$.

	\emph{Transitif:} Upamakna $\tuple{a, b} \Intequiv \tuple{c, d}\Intequiv \tuple{m, n}$. Mula $a + d = c + b$ lan $c + n = m + d$. Dadi $a + d + c + n = c + b + m + d$, mula $a + n = m + b$. Dadi $\tuple{a, b } \Intequiv \tuple{m, n}$.
\end{proof}

Saiki relasi ekuivalensi iki bisa digunakake kanggo njupuk kelas-kelas ekuivalensi:

\begin{defn}
		Wilangan wutuh yaiku kelas-kelas ekuivalensi pasangan urut wilangan natural miturut $\Intequiv$; yaiku $\Int = \equivclass{\Nat^2}{\Intequiv}$.
\end{defn}

Definisi stipulatif iki bisa nuwuhake akeh reaksi \emph{filosofis} kang beda. Nanging sadurunge ngrembug reaksi kasebut, prayoga yen kita nerusake dhisik sawetara prakara teknis.

Sawise netepake apa iku wilangan wutuh, kita kudu netepake fungsi lan relasi dhasar ing wilangan kasebut. Ayo nulis $\equivrep{m, n}{\Intequiv}$ kanggo kelas ekuivalensi miturut $\Intequiv$ kang duwe $\tuple{m, n}$ minangka !!a{element}.\footnote{Cathetan: yen nggunakake notasi kang dikenalake ing \olref[sfr][rel][eqv]{def:equivalenceclass}, kita bakal nulis $\equivrep{\tuple{m,n}}{\Intequiv}$ kanggo perkara kang padha. Nanging notasi iku rada luwih angel diwaca.} Tegese:
	$$\equivrep{m, n}{\Intequiv} = \Setabs{\tuple{a, b} \in \Nat^2}{\tuple{a, b}\Intequiv \tuple{m, n}}$$
Saiki kita menehi sawetara definisi:
	\begin{align*}
		\equivrep{a, b}{\Intequiv} + \equivrep{c, d}{\Intequiv} &= \equivrep{a + c, b + d}{\Intequiv}\\
		\equivrep{a, b}{\Intequiv} \times \equivrep{c, d}{\Intequiv} &= \equivrep{a c + b  d, a  d + b c}{\Intequiv}\\
		\equivrep{a, b}{\Intequiv} \leq \equivrep{c, d}{\Intequiv} &\text{ yen lan mung yen }a + d \leq b + c
	\end{align*}
(Kaya kang lumrah, "$ab$" digunakake kanggo ngganti "$(a \times b)$", supaya aksioma luwih gampang diwaca.) Saiki kita kudu mesthekake yen definisi-definisi iki tumindak kaya kang \emph{kudune}. Njlentrehake tegese iki lan mriksa kabeh kanthi rinci rada ngrepotake; rinciane dipindhah menyang \olref[check]{sec}. Nanging asil ringkese: kabeh lumaku kanthi becik!

Isih ana siji prakara pungkasan. Kita wis mbangun wilangan wutuh nganggo
wilangan natural. Nanging iki ateges wilangan natural \emph{dudu
wilangan wutuh kasebut dhewe}. Kita bakal mbahas teges filosofise
ing \olref[ref]{sec}. Nanging saka sisih teknis, kita mbutuhake cara
kanggo nganggep wilangan natural \emph{minangka} wilangan wutuh.
Gagasane gampang: kanggo saben $n \in \Nat$, kita netepake
$n_\Int = \equivrep{n, 0}{\Intequiv}$. Kita kudu mesthekake yen
definisi iki tumindak kanthi becik, yaiku kanggo saben $m, n
\in \Nat$
\begin{align*}
	(m + n)_\Int &= m_\Int + n_\Int\\
	(m \times n)_\Int &= m_\Int \times n_\Int\\
	m \leq n &\liff m_\Int \leq n_\Int
\end{align*}
Nanging kabeh iki cukup gampang. Upamane, kanggo nuduhake yen
sarat kapindho lumaku, kita bisa nggunakake tumindake wilangan natural
lan mikir kaya mangkene:
%\begin{proof}
%	Nganggo tumindake wilangan natural, cetha:
	\begin{align*}
%		(m + n)_\Int  = \equivrep{m + n, 0}{\Intequiv} = \equivrep{m + n, 0 + 0}{\Intequiv} = \equivrep{m, 0}{\Intequiv} + \equivrep{n, 0}{\Intequiv} = m_\Int + n_\Int\\
		(m \times n)_\Int  &= \equivrep{m \times n, 0}{\Intequiv} \\
		&= \equivrep{m \times n + 0 \times 0, m \times 0 + 0 \times n}{\Intequiv} \\
		&= \equivrep{m, 0}{\Intequiv} \times \equivrep{n, 0}{\Intequiv} \\
		&= m_\Int \times n_\Int
	\end{align*}
%\end{proof}
Kita pasrahake minangka gladhen kanggo mesthekake yen rong sarat liyane lumaku.
\begin{prob}
	Tuduhna yen $(m + n)_\Int = m_\Int + n_\Int$ lan $m \leq n \liff m_\Int \leq n_\Int$, kanggo saben $m, n \in \Nat$.
\end{prob}
\end{document}
